\section{Stripmap equivalent CSLC}
\label{sec:strip_eq_cslc}

\subsection{Product Description}
Managing SAR acquisitions in burst mode is more complex than in standard stripmap mode. Burst-mode acquisitions, such as those from TOPSAR or ScanSAR, introduce additional phase modulation due to the scanning nature of the acquisition. As a result, SAR products generated from burst-mode data cannot be directly mosaicked, since the overlapping regions between bursts exhibit phase discontinuities. To produce a coherent image, these effects must be corrected through processes such as resampling, interpolation, or sub-band extraction.

To reduce the complexity for the end user, each burst’s signal can be demodulated and the bursts subsequently mosaicked. This results in a single, continuous image that simulates a Stripmap acquisition.

The resulting Stripmap-equivalent (\acs{stripeq}) \acs{slc} or \acs{cslc} image is a full-resolution SAR product that preserves the resolution of the original input data. It seamlessly merges multiple bursts acquired in wide-swath modes like TOPSAR or ScanSAR into a unified image, maintaining the spatial consistency and phase continuity characteristic of Stripmap mode. The output retains the full information content of the original burst-mode data.

\subsubsection{Input and Auxiliary Data}
The following input and auxiliary data are considered for product generation:

\begin{longtblr}[
	label = tab:in_aux_stripeq_products,
	entry = none,
	caption = {\acs{stripeq} input and auxiliary data},
	]{
		width = \linewidth,
		colspec = {Q[l,20]Q[l,60]Q[c,20]},
		hlines,
		hline{1-Z} = {-}{0.05em},
	}
	\textbf{Name}   & \textbf{Description}   & \textbf{Mandatory (M) / Optional (O)} \\
    \acs{slc} or \acs{cslc} & L1 Single Look Complex image or co-registered L1 Single Look Complex image and annotations & M \\
    Precise Orbits & \acs{slc} precise orbits or \acs{cslc} precise reference orbits. & M \\
\end{longtblr}

\subsubsection{Main Variables}
The following tables list the main variables associated with the \acs{stripeq} products.

\begin{longtblr}[
	label = tab:var_stripeq_products,
	entry = none,
	caption = {\acs{stripeq} main variables},
	]{
		width = \linewidth,
		colspec = {Q[l,15]Q[l,52]Q[l,20]Q[l,13]},
		hlines,
		hline{1-Z} = {-}{0.05em},
	}
	\textbf{Name} & \textbf{Description} & \textbf{Dimensions} & \textbf{Data Type} \\
	\acs{slc}/\acs{cslc} & Single look complex or Co-registered Single look complex data: 
    
    - the amplitude is in \gls{beta0} convention, obtained by the calibration step using the Look-up-Table from the \acs{slc}/\acs{cslc} annotations;
    
    - the phase corresponds to the interpolated and resampled \acs{slc}/\acs{cslc} phase & polarization, acquisition\_date, azimuth\_time, range\_time & complex64 \\
    phase\_model & Topographic phase model, in rad, removed from the original image phase & acquisition\_date, azimuth\_time, range\_time & float32 \\
\end{longtblr}

The table below describes the dimensions/coordinates associated with the \acs{cslc} and \acs{gslc} main variables.

\begin{longtblr}[
	label = tab:var_stripeq_coordinates,
	entry = none,
	caption = {\acs{stripeq} main variables coordinates/dimensions},
	]{
		width = \linewidth,
		colspec = {Q[l,20]Q[l,62]Q[l,18]},
		hlines,
		hline{1-Z} = {-}{0.05em},
	}
	\textbf{Name} & \textbf{Description} & \textbf{Data Type} \\
	polarization & Acquired polarizaions & string \\
    acquisition\_date & Image acquisition dates, formatted as YYYY-MM-DD & string \\
    azimuth\_time & Time along satellite track direction & datetime64 \\
    range\_time & Time across satellite track direction & float64 \\
\end{longtblr}

\subsubsection{Main Metadata}
The following tables list the main metadata and per-pixel metadata associated with the \acs{stripeq} products.

\begin{longtblr}[
	label = tab:meta_stripeq_products,
	entry = none,
	caption = {\acs{stripeq} main metadata and per-pixel metadata},
	]{
		width = \linewidth,
		colspec = {Q[l,17]Q[l,38]Q[l,20]Q[l,13]Q[c,12]},
		hlines,
		hline{1-Z} = {-}{0.05em},
	}
	\textbf{Name} & \textbf{Description} & \textbf{Dimensions} & \textbf{Data Type} & \textbf{Per-pixel Metadata}\\
	input\_slc/cslc & Link to input \acs{slc}/\acs{cslc} & - & string & \\
    orbit\_files & Orbital file path list used for each of the images belonging to the \acs{cslc} stack & acquisition\_date & string & \\
    orbit\_coeffs & Polynomial coefficients describing the orbit associated with each input image belonging to the \acs{cslc} stack & acquisition\_date, axis, degree & float64 & \\
    orbit\_reference\
    
    \_time & Reference time associated with the orbit polynomial coefficients for each input image belonging to the \acs{cslc} stack & acquisition\_date & float64 & \\
    orbit\_type & Type of orbital data used for each of the images belonging to the \acs{cslc} stack: ETAD, POE, POD, RES & acquisition\_date & string & \\
    beta0 & LUT of the beta0 coefficients obtained from each input \acs{slc} annotation belonging to the \acs{cslc} stack & acquisition\_date, azimuth\_time, slant\_range\_time & float32 & \\
    gamma0 & LUT of the gamma0 coefficients obtained from each input \acs{slc} annotation belonging to the \acs{cslc} stack & acquisition\_date, azimuth\_time, slant\_range\_time & float32 & \\
    rcs & LUT of the Radar Cross Section (RCS) coefficients obtained from each input \acs{slc} annotation 
    belonging to the \acs{cslc} stack & acquisition\_date, azimuth\_time, slant\_range\_time & float32 & \\
    noise & LUT of thermal noise correction obtained from each input \acs{slc} annotation belonging to the \acs{cslc} stack & acquisition\_date, azimuth\_time, slant\_range\_time & float32 & \\
    phase\_deramp & The de-ramping phase (image or model) applied to each burst to obtain a continuous phase in the Stripmap Equivalent \acs{slc}/\acs{cslc}. & float32 & \\
    source\_data
    
    \_attributes & Dictionary containing the attributes associated with each image belonging to the \acs{cslc} stack & - & Dictionary & \\
    bursts\_list & A list of the input bursts & - & list of int32 & \\
    x\_sample & Pixel spacing in x direction & - & float32 & \\
    y\_sample & Pixel spacing in y direction & - & float32 & \\
\end{longtblr}

\subsubsection{Reference System Tiling and Slicing}
\label{sec:stripeq_reference_system_tiling}
\autoref{tab:stripeq_ref_system_slicing} reports the reference system, the slicing, and grid for the \acs{stripeq} products.


\begin{longtblr}[
	label=tab:stripeq_ref_system_slicing,
	entry = none,
	caption = {Reference system and Slicing for \acs{stripeq}},
	]{
		width=\linewidth,
		colspec = {Q[l,15]Q[l,40]Q[l,15]Q[l,30]},
		%vline{2-5} = {-}{},
		%vline{5-5} = {-}{},
		hlines,
		hline{1-Z} = {-}{0.05em},
		rowhead=1,
	}
	\textbf{Product} & \textbf{Reference System} & \textbf{Slicing} & \textbf{Spacing}\\
	\acs{stripeq} & Radar: azimuth/slant-range & Burst & Same of the input data \\   
\end{longtblr}

For the \acs{stripeq} product, it is planned to use the same reference radar grid and the same slicing for the ROSE-L and Sentinel-1 missions, to ensure interoperability between the products generated with the data of both missions.

\subsection{Product Requirements}

\subsubsection{Accuracy}

Given that the primary objective of this \acs{ard} product is to streamline the handling of complex burst-mode SAR acquisitions—such as those acquired in \acs{topsar} or ScanSAR modes—while preserving interoperability with standard \acs{stripeq} processing chains, it is critical to ensure that the product remains suitable for conventional amplitude- and phase-based analyses. These include, but are not limited to, amplitude image generation, interferogram computation, and coherence estimation.

Accordingly, the quality indicators have been defined to evaluate the continuity and consistency of derived products obtained from co-registered \acs{stripeq} image pairs. In particular, the assessment focuses on the continuity  of the interferometric phase and coherence maps generated from a pair of co-registered \acs{stripeq} images.

\begin{longtblr}[
	label = tab:stripeq_accuracy,
	entry = none,
	caption = {Accuracy Requirements for \acs{stripeq} product},
]{
	width = \linewidth,
	colspec = {Q[l,30] Q[l,60] Q[c,20]},
	hlines,
	hline{1-Z} = {-}{0.05em},
	rowhead = 1,
}

\textbf{Metric} & \textbf{Description} & \textbf{Requirement} \\

\textbf{Interferometric coherence discontinuity} & Maximum interferometric phase discontinuity between two adjacent bursts (only in the azimuth direction) & $\leq 0.01$ \\

\textbf{Interferometric phase discontinuity (deg)} & Maximum interferometric phase discontinuity between two adjacent bursts (in both range and azimuth directions) in case of absence of motion in the azimuth direction & $\leq 5$~deg \\

\end{longtblr}

\begin{comment}
\begin{itemize}
    \item IRF degradation wrt SLC (tool SCT) to be evaluated on few S1 \& SAOCOM data
    \item PSLR degradation wrt SLC (tool SCT) to be evaluated on few S1 \& SAOCOM data 
    \item ... degradation wrt SLC (tool SCT) to be evaluated on few S1 \& SAOCOM data
    \item geolocation accuracy (tool SCT) in correspondence of CRS
    \item Spectral information preservation (doppler difference wrt input data)
    \item Coregistration accuracy (rg, az) (theoretical formulas, CRS?) as a function of Ancillary data / processing
    \item phase continuity between bursts (theoretical analysis) ???? as a function of Ancillary data / processing
\end{itemize}    
\end{comment}

\subsubsection{Timeliness}
The timeliness is comparable to that of \acs{cslc} generation (see \textsection~\ref{sec:cslc_timeliness}). The de-ramping and de-bursting operations, necessary for generating this product, do not significantly impact timeliness.

\subsection{Usability, Integrability and Interoperability}
Usability and integrability of this product are the same as those collected for \acs{cslc} (\textsection~\ref{sec:usability_integrability_CSLC}), of which this novel product is an extension. It does not require processing related to burst fusion/merging.
